%%%%%%%%%%%%%%%%%%%%%%%%%%%%%%%%%%%%%%%%%%%%%%%%%%%%%%%%%%%%
% ABSTRACT
%%%%%%%%%%%%%%%%%%%%%%%%%%%%%%%%%%%%%%%%%%%%%%%%%%%%%%%%%%%%

\chapter*{}
\addcontentsline{toc}{chapter}{Abstract}

\begin{center}
{\Huge\bfseries Abstract}
\end{center}

\vspace{1cm}

Active droplets are self-propelled systems that move by interacting
with their surrounding chemical environment. As they move, they leave
behind a chemical trail that can influence their future motion,
leading to memory-dependent behaviour.

\medskip

In this work, droplet trajectories were extracted from experimental
video recordings and analysed using statistical measures such as the
mean squared displacement (MSD), velocity autocorrelation function
(VACF), and orientation correlation function. The results show that
the droplets exhibit persistent motion at short times and gradually
approach diffusive behaviour at longer times. The MSD indicates
superdiffusive transport, while the correlation functions reveal
finite velocity and directional memory.

\medskip

To understand these observations, a self-repelling random-walk
framework and a coupled partial differential equation--stochastic
differential equation (PDE--SDE) model were studied. In these models,
the droplet interacts with its own chemical trail, which stores
information about past motion and produces non-Markovian dynamics.

\medskip

The agreement between experimental observations and model predictions
suggests that the chemical trail acts as a memory field that
influences future motion. These results demonstrate that memory
effects play an important role in the dynamics of active droplets and
provide a useful framework for understanding non-Markovian transport
in active matter systems.

\vspace{1cm}

\noindent
\textbf{Keywords:} Active droplets, active matter, memory effects,
non-Markovian dynamics, mean squared displacement, velocity
autocorrelation function, PDE--SDE model.

\newpage
